\begin{lemma}[Symmetry of the second weak derivatives]
Let $n\ge1$ be an integer and $q\ge1$ real, let $\Omega\subseteq\R^n$ be open,
let $u\in W^{1,q}_{\mathrm{loc}}(\Omega)$ with weak gradient $g$
(\noderef{MemW1pLocOn}), and suppose that for each $i\in\{1,\dots,n\}$ the
component $g^i$ again lies in $W^{1,q}_{\mathrm{loc}}(\Omega)$ with weak
gradient $Dg^i$ (\noderef{MemW1pLocOn}).  Then for all
$j,k\in\{1,\dots,n\}$,
\[
(Dg^k)_j=(Dg^j)_k\qquad\text{almost everywhere on }\Omega ,
\]
that is, the second weak derivatives $D_jD_ku$ and $D_kD_ju$ of $u$ agree almost
everywhere on $\Omega$.
\end{lemma}
\begin{proof}
Fix $j,k\in\{1,\dots,n\}$ and put
\[
w=(Dg^k)_j-(Dg^j)_k ,
\]
a real-valued function on $\R^n$.  We show that $\int_{\R^n}\varphi\,w\,dx=0$
for every $\varphi\in C^\infty(\R^n)$ with compact support and
$\operatorname{supp}\varphi\subseteq\Omega$, and that $w$ is locally integrable
on $\Omega$; the fundamental lemma of the calculus of variations (du
Bois--Reymond, Preamble) then gives $w=0$ almost everywhere on $\Omega$, which
is the assertion.

\emph{Step 1: local integrability, and the finiteness of every integral below.}
Let $K\subseteq\Omega$ be compact.  By \noderef{MemW1pLocOn} applied to $u$ with
weak gradient $g$ we have $u\in L^q(K)$ and $|g|\in L^q(K)$, both almost
everywhere strongly measurable there; by \noderef{MemW1pLocOn} applied to the
component $g^i$ with weak gradient $Dg^i$ we have $g^i\in L^q(K)$ and
$|Dg^i|\in L^q(K)$ for each $i$.  Since $K$ is compact, $|K|<\infty$, and
$q\ge1$, so H\"older on a set of finite measure (Preamble) gives
$L^q(K)\subseteq L^1(K)$; hence $u$, each $g^i$ and each $|Dg^i|$ lie in
$L^1(K)$.  Each component satisfies $|(Dg^i)_m|\le|Dg^i|$ pointwise, so each
$(Dg^i)_m$ lies in $L^1(K)$ as well, by domination (Preamble).  Consequently
$w=(Dg^k)_j-(Dg^j)_k$ is measurable and lies in $L^1(K)$ for every compact
$K\subseteq\Omega$; since $\Omega$ is open, every point of $\Omega$ has a
compact neighbourhood contained in $\Omega$, so $w$ is locally integrable on
$\Omega$.

\emph{Step 2: the test functions.}  Fix $\varphi\in C^\infty(\R^n)$ with compact
support and $\operatorname{supp}\varphi\subseteq\Omega$, so that $\varphi$
belongs to the test class of \noderef{HasWeakGradientOn} for $\Omega$.  For each
index $m$ the partial derivative $D_m\varphi$ is again in $C^\infty(\R^n)$, and
$\{D_m\varphi\ne0\}\subseteq\operatorname{supp}\varphi$ because $\varphi$
vanishes identically on the open complement of its support, so
$\operatorname{supp}D_m\varphi\subseteq\operatorname{supp}\varphi$ is compact
and contained in $\Omega$.  Hence $D_m\varphi$ is admissible in
\noderef{HasWeakGradientOn} for $\Omega$ too.  The same applies to the second
derivatives $D_mD_{m'}\varphi$.  All of $\varphi$, $D_m\varphi$ and
$D_mD_{m'}\varphi$ are continuous with compact support, hence bounded on
$\R^n$, and all vanish off the compact set
$K_0=\operatorname{supp}\varphi\subseteq\Omega$.  By Step 1, every integral
written below is therefore the integral over $K_0$ of an $L^1(K_0)$ function
times a bounded function, hence finite (domination, Preamble); and since each
integrand vanishes off $K_0\subseteq\Omega$, its integral over $\Omega$ and its
integral over $\R^n$ coincide.

\emph{Step 3: two applications of the weak-gradient identity.}  The defining
identity of \noderef{HasWeakGradientOn}, which the pair $(g^k,Dg^k)$ satisfies
on $\Omega$ by \noderef{MemW1pLocOn}, applied to the test function $\varphi$ and
the index $j$, reads
\[
\int_\Omega g^k\,D_j\varphi\,dx=-\int_\Omega (Dg^k)_j\,\varphi\,dx .
\tag{1}
\]
The defining identity of \noderef{HasWeakGradientOn}, which the pair $(u,g)$
satisfies on $\Omega$ by \noderef{MemW1pLocOn}, applied to the test function
$D_j\varphi$ --- admissible by Step 2 --- and the index $k$, reads
\[
\int_\Omega u\,D_k\bigl(D_j\varphi\bigr)dx=-\int_\Omega g^k\,D_j\varphi\,dx ,
\tag{2}
\]
the right-hand side being $-\int_\Omega g^k\,D_j\varphi$ because the $k$-th
component of $g$ is $g^k$.  Adding (1) and (2),
\[
\int_\Omega (Dg^k)_j\,\varphi\,dx
=-\int_\Omega g^k\,D_j\varphi\,dx
=\int_\Omega u\,D_kD_j\varphi\,dx .
\tag{3}
\]

\emph{Step 4: the right-hand side of (3) is symmetric in $(j,k)$.}  The function
$\varphi$ is $C^\infty$, hence differentiable at every point of $\R^n$, and its
derivative $D\varphi$ is again differentiable at every point.  By the symmetry
of the second derivative (Clairaut, Preamble), which requires exactly Fr\'echet
twice differentiability and no more, the bilinear map $D^2\varphi(x)$ is
symmetric for every $x$:
$D^2\varphi(x)[v][w]=D^2\varphi(x)[w][v]$ for all $v,w\in\R^n$.  Taking
$v=e_k$, $w=e_j$ and using $D_m\varphi=D\varphi[e_m]$, so that
$D_kD_j\varphi(x)=D^2\varphi(x)[e_k][e_j]$ and
$D_jD_k\varphi(x)=D^2\varphi(x)[e_j][e_k]$, we get
\[
D_kD_j\varphi=D_jD_k\varphi\qquad\text{everywhere on }\R^n .
\]
Hence the right-hand side of (3) is unchanged when $j$ and $k$ are exchanged.
Exchanging $j$ and $k$ in (3) and comparing therefore gives
\[
\int_\Omega (Dg^k)_j\,\varphi\,dx=\int_\Omega (Dg^j)_k\,\varphi\,dx ,
\]
both integrals being finite by Step 2, so that they may be subtracted:
\[
\int_{\R^n}\varphi\,w\,dx=\int_\Omega\bigl[(Dg^k)_j-(Dg^j)_k\bigr]\varphi\,dx=0 ,
\]
the first equality because the integrand vanishes off
$\operatorname{supp}\varphi\subseteq\Omega$ (Step 2).

\emph{Conclusion.}  The function $w$ is locally integrable on the open set
$\Omega$ by Step 1, and $\int_{\R^n}\varphi\,w=0$ for every $\varphi$ that is
$C^\infty$ with compact support contained in $\Omega$ by Step 4.  The
fundamental lemma of the calculus of variations (Preamble), whose test class is
exactly this one, yields $w=0$ almost everywhere on $\Omega$, that is
$(Dg^k)_j=(Dg^j)_k$ almost everywhere on $\Omega$.  As $j$ and $k$ were
arbitrary, this is the assertion.
\end{proof}
